% 87-01 ③(87쪽) — 보기 그래프: $x=-2$ 에서 극대, $x=1$ 에서 기울기 0 인 위로 볼록한 사차함수 개형. 그림 자체가 보기라 TikZ 로 다시 그림.
% 라벨은 원문 그대로 -2, O, 1, x, y, y=f(x) 와 점선 한 줄만.
\begin{tikzpicture}[baseline=(current bounding box.center), x=0.40cm, y=0.40cm, line width=0.5pt]
  \path (-3.6,-2.2) rectangle (3.45,3.6);
  \draw[->, >=stealth, line width=0.7pt] (-3.5,0) -- (2.7,0) node[below=1pt, font=\footnotesize] {$x$};
  \draw[->, >=stealth, line width=0.7pt] (0,-1.9) -- (0,3.3) node[left=1pt, font=\footnotesize] {$y$};
  \draw[densely dotted, line width=0.4pt] (-2,0) -- (-2,2.62);
  \draw[line width=0.6pt, smooth] plot coordinates {(-3.3,-1.7) (-3.15,-0.9) (-3,-0.15) (-2.85,0.6) (-2.7,1.25) (-2.5,1.95) (-2.3,2.4) (-2.15,2.58) (-2,2.62) (-1.85,2.56) (-1.65,2.35) (-1.4,1.95) (-1.15,1.55) (-0.85,1.15) (-0.5,0.78) (-0.1,0.45) (0.3,0.23) (0.65,0.12) (1,0.08) (1.3,0.05) (1.5,-0.05) (1.7,-0.3) (1.9,-0.75) (2.05,-1.25) (2.15,-1.7)};
  \node[below=1pt, font=\footnotesize] at (-2,0) {$-2$};
  \node[below left=1pt and 1pt, font=\footnotesize] at (0,0) {$\pt{O}$};
  \node[below=1pt, font=\footnotesize] at (1,0) {$1$};
  \node[anchor=west, font=\footnotesize] at (0.28,1.0) {$y=f(x)$};
\end{tikzpicture}
